% ECONOMICS_PROBLEM: {"id": "D86-132", "title": "Approximation frontier for randomized inspection contracts", "jel_code": "D86", "source_id": "OP-0597"}
% !TeX program = xelatex
\documentclass[11pt,a4paper]{article}
\usepackage[margin=23mm]{geometry}
\usepackage{fontspec}
\setmainfont{Latin Modern Roman}
\usepackage{amsmath,amssymb,mathtools}
\usepackage{xcolor,enumitem,booktabs,longtable,array,xurl,hyperref}
\definecolor{muted}{HTML}{53636C}
\hypersetup{colorlinks=true,urlcolor=blue,linkcolor=blue}
\setlength{\parindent}{0pt}
\setlength{\parskip}{5pt}
\newcommand{\R}{\mathbb R}
\newcommand{\E}{\mathbb E}
\newcommand{\N}{\mathbb N}
\newcommand{\ind}{\mathbf 1}
\newcommand{\Var}{\operatorname{Var}}
\newcommand{\Cov}{\operatorname{Cov}}
\newcommand{\supp}{\operatorname{supp}}
\DeclareMathOperator*{\argmin}{arg\,min}
\DeclareMathOperator*{\argmax}{arg\,max}
\newcommand{\entrymeta}[3]{{\sffamily\small\color{muted}\textbf{\texttt{#1}}\quad|\quad #2\par #3\par}}
\newcommand{\Problemref}[1]{\texttt{#1}}
\newcommand{\JELref}[2]{\texttt{#2} (JEL \texttt{#1})}
\begin{document}
\section*{Approximation frontier for randomized inspection contracts}
\textbf{Problem ID:} \texttt{D86-132}\quad \textbf{JEL:} \texttt{D86}\par
% BEGIN ECONOMICS STATEMENT
\label{problem:OP-0597}
% ECONOMICS_PROBLEM: {"local_id": "ECON-012", "title": "Approximation frontier for randomized inspection contracts", "kind": "P", "audit_date": "2026-10-08", "source_identity_checked": true, "agenda_or_question_support_checked": true, "source_role": "see per-entry audit and locators", "current_open_certified": false}


\label{audited:ECON-012}

{\sffamily\small\color{muted}\textbf{ECON-012} | P | Explicit published open question, at source date\par}

\subsection*{Definitions and formal target}

Let \(A\) be \(n\) actions with known costs \(c_a\geq0\) and success probabilities \(f_a\in[0,1]\), including a null action of cost zero. Success pays the principal one. A committed contract recommends \(a\), pays \(\alpha\in[0,1]\) upon success if no deviation is detected, and independently draws an undisclosed inspection set \(S\subseteq A\) with law \(p\). A deviation to \(b\ne a\) is detected when \(a\in S\) or \(b\in S\).

The normalized monotone inspection cost \(h:2^A\to\mathbb R_{\geq0}\) has \(h(\varnothing)=0\). XOS means \(h(S)=\max_{\ell\in K}\sum_{a\in S}w_{\ell a}\) for a finite family with \(w_{\ell a}\geq0\); the family is not assumed given explicitly. Subadditive means \(h(S\cup T)\leq h(S)+h(T)\). Submodular means \(h(S)+h(T)\geq h(S\cup T)+h(S\cap T)\).

Agent utilities and the compliant principal payoff are
\[
 U_a=\alpha f_a-c_a,\quad
 U_b=\alpha f_b\Pr_{S\sim p}(S\cap\{a,b\}=\varnothing)-c_b,\quad
 \Pi=(1-\alpha)f_a-\mathbb E_p h(S).
\]
Require \(U_a\geq U_b\) for every \(b\), with ties selecting \(a\); let \(\Pi^\star\) be the maximum over such contracts. Under the source's exact value/demand oracle and arithmetic conventions, does XOS admit a PTAS, \(\Pi\geq(1-\varepsilon)\Pi^\star\), for every fixed \(\varepsilon\in(0,1)\)? Does the subadditive class admit a constant \(C\geq1\) with \(\Pi\geq\Pi^\star/C\)?

\subsection*{Clarifications and required assumptions}

A cost-zero action with zero payment and no inspection provides nonnegative principal payoff, so \(\Pi^\star\geq0\); its success probability need not be zero. Inspection is independent of success and the agent acts before observing the draw. The finite distribution over \(2^A\) defines the mathematical feasible set; a polynomial algorithm must output an implementable distribution, with a declared support/sampling representation. A value query returns \(h(S)\); a demand query at nonnegative prices \(q\) maximizes \(h(S)-\sum_{a\in S}q_a\). The source treats numeric descriptions as \(O(1)\). A rational-input polynomial-bit guarantee is a separately specified extension, not silently supplied by that convention.

\subsection*{Variants and changes of scope}

\textbf{XOS PTAS.} For each fixed $\varepsilon\in(0,1)$, require polynomial time in the number of actions under the declared arithmetic/oracle model, with $\Pi\geq(1-\varepsilon)\Pi^\star$. \textbf{Subadditive approximation.} Require $\Pi\geq\Pi^\star/C$ for an absolute $C\geq1$, rather than assuming a PTAS. \textbf{Oracle strength.} The source permits value and demand access; a demand query maximizes $h(S)-\sum_{g\in S}q_g$. Value-only access is a stronger algorithmic request and does not automatically inherit the same open status. \textbf{Observation.} Fully observable inspection, partial observation, and deterministic inspection are different problems.

\subsection*{Audit conclusion and source relationship}

The inspection payoff formula matches Section 2. Section 8 explicitly discusses the XOS PTAS and subadditive constant-factor questions. The draft had silently weakened oracle access to value-only; the main question now uses the source's value/demand access model. The payoff guarantee is multiplicative, not an additive approximation.

\subsection*{References and exact source roles}

{\footnotesize\raggedright
\par\noindent\textbf{[S1]} Tomer Ezra, Stefano Leonardi, and Matteo Russo (2026). \emph{Contract Design Beyond Hidden-Actions}. SODA 2026, 2113--2133; DOI: 10.1137/1.9781611978971.77. Checked author version: arXiv:2402.16553v2, 26 November 2025.\par \textit{Verified locator / source role:} Section 2 (model and value/demand oracles); Section 8, Discussion (approximation questions). The draft incorrectly called Section 8 a conclusion.\par \url{https://arxiv.org/html/2402.16553v2}\par\smallskip
}

\subsection*{Common conventions: Additional-source status and mathematical conventions}

\label{audited:conventions}

Of the 100 supplied ECON entries, 65 distinct problems appear here; 32 close
matches refine earlier entries, and three duplicate questions map to existing
entries. Appendix~\ref{app:audited-crosswalk} lists every destination.
The attachment's P/R/S/U distinctions and source audits are retained. Its
precise mathematical formalizations are editorial unless the individual audit
says otherwise. Candidate subsidy targets ECON-004--006 retain their U flags;
the resolved approval-core question ECON-007 updates OP-0202 above.
The following supplied conventions govern both this part and the attributed
ECON refinements elsewhere in the catalogue, subject to local restrictions.

\textbf{Parameterized questions.} A problem family lists inputs that must be fixed before an instance is evaluated: population, horizon, policies, information, data law, model class and welfare weights. These are required choices, not quantities the citations are assumed to specify. Undefined applications should not be treated as identified estimands.

\textbf{Indices and integrability.} Sets of agents, firms and regions are finite unless a population measure is explicitly used. \(i,j\) index units, \(r\) regions, \(t\) dates and \(h\) horizons. \(\mathbb E\) and \(\Pr\) refer to the declared population and law; \(\mathbf1\{A\}\) is an indicator. Every displayed expectation and discounted sum needs existence in the stated domain. Infinite horizons use a discount in \((0,1)\) and a convergence condition; finite horizons may use discount one.

\textbf{Optimization and existence.} A displayed \(\max\), \(\min\), or \(\arg\max\) is conditional on attainment. For finite feasible menus this is immediate; general spaces need continuity and compactness/coercivity or another existence argument. Without attainment, the target is the supremum/infimum and arbitrarily near-optimal feasible choices. \(\inf\varnothing=+\infty\). Empty feasibility and an unidentified target are different issues.

\textbf{Potential outcomes.} \(Y_i(z)\) is the outcome under a specified intervention, not the observed conditional mean among people choosing \(z\). In binary comparisons \(\Delta Y_i=Y_i(1)-Y_i(0)\). Specify assignment, treatment versions, compliance, information and observation horizon. Interference is allowed under a full policy vector; compressed exposure notation needs a justified mapping. No interference, consistency and overlap are substantive assumptions, not automatic consequences of notation.

\textbf{Populations and observation.} Fix baseline eligibility and population weights. Conditioning on post-policy employment, survival, relocation or program uptake can change the population being compared. Death is not ordinary loss to follow-up; outcomes truncated by death need a composite convention, joint survival target, or explicitly defined survivor stratum. Fertility-changing interventions need a population functional rather than an unexplained fixed descendant cohort. Measurement and missingness models must be supplied.

\textbf{Identification.} A structural model \(m\in\mathcal M\) implies observable law \(P_m\) and target \(\theta(m)\). Identification means
\[
 P_{m_1}=P_{m_2}\ \Longrightarrow\ \theta(m_1)=\theta(m_2)
 \quad\text{for all }m_1,m_2\in\mathcal M .
\]
Without identification, the target is
\[
 \Theta(P;\mathcal M)=\{\theta(m):m\in\mathcal M,\ P_m=P\}.
\]
Writing \(\theta(P)\) before establishing identification can hide incompatible latent models. Confidence sets additionally need a sampling law, confidence level, asymptotic sequence and uniformity class. Point identification, coverage and informative precision are separate properties.

\textbf{Mediation.} Total effects of randomized assignment are distinct from cross-world natural indirect effects. Belief/effort-channel effects require a meaningful mediator intervention and additional measurement, exclusion or mediation assumptions. Randomization of the initial policy alone does not supply those assumptions. Report bounds or interventional alternatives when the stronger target is unsupported.

\textbf{Equilibrium.} A counterfactual \(W(z)\), \(p(z)\), or \(\mu_t^z\) requires individual optimization, feasibility, government/resource budgets, market clearing and state-transition laws. State equilibrium existence and selection, or report ranges over compatible equilibria. Initial-state integration includes subsequent endogenous transitions; current-state integration must use the policy-dependent distribution. Reduced-form invariance after policy is not assumed.

\textbf{Welfare accounts.} Scores, utility, health, dollars and ecological quantities require explicit conversion before addition. Fees, taxes, subsidies, wages and extortion payments are transfers between parties; distributional weights can make them welfare relevant, but their face value is not automatically a resource loss. State ownership, geographic boundaries, foreign-income weights and fiscal finance. Do not add profits to household welfare already including dividends, or count land capitalization and the underlying benefit twice. A benefit-cost expression must distinguish gross benefits from net welfare and subtract every real cost exactly once.

\textbf{Derivatives and risk.} Log derivatives require positive arguments; local responses require admissible perturbations and specified equilibrium branches. Differentiation under expectations needs regularity/domination, and boundaries may allow only one-sided derivatives. Risk uses a specified law; ambiguity uses a declared nonempty law set on a common history space. Adaptive policies must be nonanticipative. A robust criterion is an editorial choice unless attributed explicitly.

\textbf{Scope overlaps.} Allocation, collective choice, contracts, household bargaining and coordinated policy overlap economics with optimization and game theory. They are retained because they appeared in the original catalogue; no claim of a strictly game-theory-free selection is made.
% END ECONOMICS STATEMENT
\end{document}
